\documentclass[11pt]{article}
% === core_packages.tex ===
\usepackage{amsmath, amssymb, amsthm, amsfonts}
\usepackage{geometry}
\usepackage{hyperref}
\usepackage{graphicx}
\usepackage{physics}
\usepackage{booktabs}
\usepackage{listings}
\usepackage{bookmark}
\usepackage{amscd}
\usepackage{tikz-cd}
\usepackage{mathtools}
\usepackage{xspace}
\usepackage[backend=biber, style=numeric]{biblatex}
\usepackage{tcolorbox}
\tcbuselibrary{theorems}
\usepackage{parskip}
\usepackage{tikz}
\usepackage{etoolbox}
\usetikzlibrary{arrows.meta, shapes, positioning, calc, cd, shapes.geometric, backgrounds}
\usepackage[en-US]{datetime2}
\usepackage{array}


% --- Math & Ontology Macros ---
\newcommand{\entropy}{S}
\newcommand{\opponent}{opponent}
\newcommand{\aside}[1]{[\emph{\opponent}: ``#1'']}
\newcommand{\paper}[1]{[\emph{#1}: ``#1'']}
\newcommand{\iame}{Real\textsuperscript{e}}
\newcommand{\iameverything}{Real\textsuperscript{everything}}
\newcommand{\fabric}{voxel}
\newcommand{\Fabric}{Voxel}
\newcommand{\pdflink}[1]{}



% --- Global Layout Defaults ---
\geometry{letterpaper, margin=1in}
\hypersetup{
    colorlinks=true,
    linkcolor=black,
    citecolor=black,
    filecolor=magenta,
    urlcolor=blue
}
\hfuzz=5pt \vfuzz=5pt \hbadness=10000 \vbadness=10000

\author{\iame}


\usepackage{tocloft}

% --- Article Specific Theorems ---
\newtheorem{lemma}{Lemma}
\theoremstyle{definition}
\newtheorem{axiom}{Axiom}
\newtheorem{theorem}{Theorem}
\newtheorem{corollary}{Corollary}
\newtheorem{definition}{Definition}
\newtheorem{proposition}{Proposition}
\newtheorem{conjecture}{Conjecture}


\makeatletter
\let\remark\@undefined
\let\endremark\@undefined
\makeatother
\theoremstyle{remark}
\newtheorem*{remark}{Remark}

\newtcbtheorem[number within=section]{principle}{Principle}%
{colback=blue!5,colframe=blue!50!black,fonttitle=\bfseries}{principle}

% --- Article TOC Layout ---
\setlength{\cftbeforesecskip}{1.0em}
\renewcommand{\cftsecleader}{\cftdotfill{\cftdotsep}}
\renewcommand{\cftsubsecleader}{\cftdotfill{\cftdotsep}}


\addbibresource{../references.bib}

\title{Fermion Structure and Particle Classification\\
from Codec Geometry}
\author{Juha Meskanen}
\date{June 2026}

\begin{document}

\maketitle

\begin{abstract}
We prove three exact theorems about the density matrix of a fermion
in equal superposition across $n$ sites, within the codec framework
of Papers~I--VII.
Let $F = \operatorname{diag}(\rho)$ be the fermionic diagonal and
$B = \rho - \operatorname{diag}(\rho)$ the bosonic off-diagonal.
Then:
\begin{align*}
  \|F\| &= \frac{1}{\sqrt{n}} \quad \text{(Universal Fermion Norm)}, \\
  \|B\| &= \sqrt{\frac{n-1}{n}} \quad \text{(Universal Boson Norm)}, \\
  \|F\|^2 + \|B\|^2 &= 1 \quad \text{(Conservation Law)}.
\end{align*}
All three results are phase-independent and hold for arbitrary
$n \geq 1$.
The winding number $m$ of the phase pattern determines the internal
symmetry of the fermion's compression residual.

The central claim of this paper differs from the Standard Model in
a foundational respect: bosons are not fundamental particles.
They are the codec's internal record of fermion transitions, and
are never directly observed.
Every physical observation is a fermion observation.
The particle classification developed here is therefore a
\emph{fermion} classification: the pair $(n, m)$ describes the
fermion's complexity class, and what Standard Model language calls
a ``gauge boson'' is the name we give to correlations between
fermion events that would otherwise require an unexplained interaction
term.
Under this reframing, the open problems of the framework reduce
to purely fermionic questions: what distinguishes the charged lepton
generations, why fermions carry the charges they carry, and why
their trajectories produce the force laws they produce.
A combinatorial argument for a 4-bit codec unit per fermion hop is
noted but is explicitly not connected by any derivation to the
observed lepton mass gaps, which are unequal (7.69 and 4.07 bits)
and do not both round to a common unit.
\end{abstract}

\tableofcontents
\newpage

%------------------------------------------------------------
\section{Introduction}
\label{sec:intro}

The codec framework of Papers~I--VII treats the universe as a finite
static bitstring of $n \approx 184$ bits.
Fermions are localised excitations --- pixels --- encoded as
normalised complex wavefunctions.
Bosons are compression residuals: the off-diagonal entries of the
density matrix $\rho = |\psi\rangle\langle\psi|$ that remain after
the per-site fermionic diagonal is removed.

Paper~VII showed that a two-site fermion hop produces a compression
residual $B$ with $\|B\| = \sin(2\theta)/\sqrt{2}$, purely
antisymmetric structure, and a universal $-\pi/2$ phase encoding the
Born rule exactly.
That result is about the \emph{fermion's motion} --- the boson is
what the codec records when a fermion moves, not a separate particle
that mediates the motion.

The present paper asks: what is the universal structure of the
fermionic content itself, and how does it organise into the particle
spectrum we observe?

The answer has two parts.
First, three exact theorems establish the norm structure of any
$n$-site fermion, culminating in a conservation law that subsumes
all previous results.
Second, the pair $(n, m)$ --- sites and winding number --- classifies
the fermion complexity classes.
The classification is purely fermionic: no interaction terms,
coupling constants, or gauge symmetries are introduced.
What Standard Model language calls bosons are derived labels for
correlations between fermion events, not independent objects.

A note on scope: gravity was derived in Papers~IV--V at the geometric
layer, independently of the wavefunction codec.
That derivation is not revisited here.
The relationship between the geometric and codec layers is addressed
in the Discussion.

%------------------------------------------------------------
\section{Setup}
\label{sec:setup}

We follow the conventions of Paper~VII exactly.
A \emph{fermion} is a localised excitation encoded as a normalised
complex wavefunction $\psi \in \mathbb{C}^n$.
The \emph{density matrix} is $\rho = |\psi\rangle\langle\psi|$,
with components $\rho_{ij} = \psi_i\psi_j^*$.
We split $\rho$ into two orthogonal parts:
\[
  F = \operatorname{diag}(\rho), \qquad
  B = \rho - \operatorname{diag}(\rho),
\]
where $F$ is the fermionic diagonal (per-site occupancy, real and
non-negative) and $B$ is the off-diagonal matrix encoding
inter-site correlations.

An \emph{$n$-site equal superposition} is
\[
  \psi_k = \frac{1}{\sqrt{n}}\,e^{i\phi_k}, \quad k = 0,\ldots,n-1,
\]
for arbitrary real phases $\phi_k$.
The natural phase patterns on an $n$-site ring are the discrete
Fourier modes
\[
  \phi_k^{(m)} = \frac{2\pi mk}{n},
  \quad m = 0, 1, \ldots, \lfloor n/2 \rfloor,
\]
where $m$ is the \emph{winding number} --- the number of full phase
rotations around the ring.
These are the minimum-$C_s$ states on the ring, ordered by ascending
spectral complexity.

%------------------------------------------------------------
\section{Three Exact Theorems}
\label{sec:theorems}

\begin{theorem}[Universal Boson Norm]
\label{thm:boson}
For any $n$-site equal superposition with arbitrary phases,
\[
  \|B\| = \sqrt{\frac{n-1}{n}}.
\]
\end{theorem}

\begin{proof}
$\rho_{ij} = (1/n)\,e^{i(\phi_i-\phi_j)}$, so
$|B_{ij}|^2 = 1/n^2$ for all $i \neq j$.
There are $n(n-1)$ off-diagonal pairs, giving
$\|B\|^2 = n(n-1)/n^2 = (n-1)/n$. \qed
\end{proof}

\begin{theorem}[Universal Fermion Norm]
\label{thm:fermion}
For any $n$-site equal superposition with arbitrary phases,
\[
  \|F\| = \frac{1}{\sqrt{n}}.
\]
\end{theorem}

\begin{proof}
$F_k = |\psi_k|^2 = 1/n$ for all $k$,
so $\|F\|^2 = \sum_k 1/n^2 = 1/n$. \qed
\end{proof}

\begin{theorem}[Fermion-Boson Conservation Law]
\label{thm:conservation}
For any pure state $\psi$,
\[
  \|F\|^2 + \|B\|^2 = 1.
\]
\end{theorem}

\begin{proof}
For any pure state, $\rho^2 = \rho$, so
$\|\rho\|^2 = \operatorname{Tr}(\rho^2) = \operatorname{Tr}(\rho) = 1$.
Since $F$ and $B$ occupy orthogonal subspaces,
$\|\rho\|^2 = \|F\|^2 + \|B\|^2 = 1$. \qed
\end{proof}

\begin{remark}[Connection to Paper~VII]
For $\psi(\theta) = \cos\theta\,e_i + i\sin\theta\,e_j$,
$\|B(\theta)\|^2 = \sin^2(2\theta)/2$ (Paper~VII, exact).
Then $\|F(\theta)\|^2 = \cos^4\theta + \sin^4\theta$, and
\[
  \|F\|^2 + \|B\|^2
  = \cos^4\theta + \sin^4\theta + 2\sin^2\theta\cos^2\theta
  = (\cos^2\theta + \sin^2\theta)^2 = 1.
\]
Theorem~\ref{thm:conservation} subsumes the Paper~VII result as a
special case.
\end{remark}

\begin{corollary}[Pauli exclusion]
For a localised fermion $\psi = e_j$: $\|F\| = 1$, $\|B\| = 0$.
No compression residual is produced.
Double occupancy has zero off-diagonal content and is
informationally invisible.
\end{corollary}

Table~\ref{tab:norms} summarises the theorems for small $n$.

\begin{table}[h]
\centering
\begin{tabular}{ccccc}
\toprule
$n$ & $\|F\| = 1/\sqrt{n}$ & $\|B\| = \sqrt{(n-1)/n}$
    & $\|F\|^2$ & $\|B\|^2$ \\
\midrule
1 & $1$          & $0$          & $1$   & $0$   \\
2 & $1/\sqrt{2}$ & $1/\sqrt{2}$ & $1/2$ & $1/2$ \\
3 & $1/\sqrt{3}$ & $\sqrt{2/3}$ & $1/3$ & $2/3$ \\
4 & $1/2$        & $\sqrt{3}/2$ & $1/4$ & $3/4$ \\
\bottomrule
\end{tabular}
\caption{Universal fermion and boson norms.
All rows satisfy $\|F\|^2 + \|B\|^2 = 1$ exactly.}
\label{tab:norms}
\end{table}

%------------------------------------------------------------
\section{Fermion Spin Structure}
\label{sec:halfinteger}

The winding number $m$ determines the symmetry of the fermion's
compression residual.
The half-integer spin of the fermion itself is encoded separately,
in the universal phase structure established in Paper~VII.

For any single fermion hop, the off-diagonal entries of $B$ carry a
universal $-\pi/2$ phase.
Four successive hops accumulate a total phase of $4 \times (-\pi/2)
= -2\pi$, returning the system to its original state.
This is precisely the spinor double-cover structure of spin-$1/2$:
a $2\pi$ rotation returns a fermion to its state only up to a sign,
while a $4\pi$ rotation is the identity.

The conservation law therefore unifies two structures in a single
codec object:
\begin{itemize}
  \item The diagonal $F$ carries half-integer spin, inherited from
        the $-\pi/2$ universal phase of the fermion hop.
  \item The off-diagonal $B$ carries the internal symmetry class
        determined by the winding number $m$.
\end{itemize}
This is not a postulate; it follows from the imaginary encoding
$\psi = f_0 + if_1$ and the definition of $B$, as shown in Paper~VII.

%------------------------------------------------------------
\section{Fermion Complexity Classes}
\label{sec:classes}

The norms $\|F\|$ and $\|B\|$ depend only on $n$ --- the number of
sites across which the fermion is delocalised.
The internal structure of the compression residual is determined
by the winding number $m$.
Together, $(n, m)$ defines a \emph{fermion complexity class}.

The symmetry of the off-diagonal $B$ under transposition determines
whether the fermion can propagate freely or is confined:

\begin{itemize}
  \item $B = -B^\top$ (antisymmetric, $m$ odd for $n{=}2$):
        the fermion's residual has self-contained antisymmetric
        structure and can propagate as a free asymptotic state.
  \item $B = B^\top$ (symmetric, $m{=}0$):
        the residual has self-contained symmetric structure
        and propagates freely with scalar character.
  \item Mixed symmetry ($n \geq 3$, $m \geq 1$):
        the residual has no self-contained symmetry class.
        The fermion requires reference to its source configuration
        to be described --- it cannot propagate as a free asymptotic
        state.
        This is the information-theoretic basis of confinement
        (see Section~\ref{sec:confinement}).
\end{itemize}

%------------------------------------------------------------
\section{The $n{=}2$ Fermion: Leptons}
\label{sec:n2}

The $n{=}2$ sector is the simplest non-trivial fermion: a two-site
equal superposition with $\|F\| = \|B\| = 1/\sqrt{2}$.

\subsection{Charged lepton ($m=1$)}

For $\psi = (e_0 + ie_1)/\sqrt{2}$ (winding $m{=}1$), the
compression residual $B$ is purely antisymmetric with eigenvalues
$\pm 1/2$ and the universal $-\pi/2$ phase of Paper~VII.
The fermion has a self-contained antisymmetric residual and
propagates freely.
This is the charged lepton.

What Standard Model language calls a ``photon'' is the name given to
the correlation between two such fermion events: when one charged
lepton recoils, another does so in a correlated way.
The codec encodes this correlation in the off-diagonal of $\rho$.
No photon particle is required; the fermion's complexity class
already contains the full correlation structure.

\subsection{Neutral fermion ($m=0$)}

For $\psi = (e_0 + e_1)/\sqrt{2}$ (winding $m{=}0$), the
residual $B$ is purely symmetric with eigenvalues $\pm 1/2$ and
zero net phase winding.
The fermion has zero electromagnetic phase structure: it is
electrically neutral.
This is the neutrino.

Within the $n{=}2$ sector, the winding number acts as the charge
selector:
\begin{itemize}
  \item $m = 0$: zero phase winding $\Rightarrow$ neutral fermion
        (neutrino).
  \item $m = 1$: one phase winding $\Rightarrow$ charged fermion
        (electron, muon, or tau, depending on complexity level).
\end{itemize}
Electric charge is the $\mathrm{U}(1)$ winding number within the
$n{=}2$ sector.
No additional postulate is required.

%------------------------------------------------------------
\section{The $n{=}3$ Fermion: Quarks and Confinement}
\label{sec:n3}

\subsection{Fractional occupancy}

For a three-site equal superposition, each site carries probability
$|\psi_k|^2 = 1/3$.
This is the per-site occupancy of a quark in a colour triplet.
The value $1/3$ requires no postulate: it is the Born-rule
probability of finding the fermion on any one colour site.

\subsection{Colour structure ($m=1$)}

For winding $m{=}1$, the compression residual has off-diagonal
entries with equal magnitude:
\[
  |B_{ij}| = \frac{1}{3} \quad \text{for all colour pairs } i \neq j.
\]
This is exact colour symmetry, emerging from the equal-superposition
constraint and Theorem~\ref{thm:boson}.
No $\mathrm{SU}(3)$ symmetry is imposed.

What Standard Model language calls a ``gluon'' is the name given to
the correlation between quark colour changes.
The mixed symmetry of $B$ for $n{=}3$, $m{=}1$ means this
correlation has no self-contained symmetry class --- see
Section~\ref{sec:confinement}.

\subsection{Scalar sector ($m=0$)}

For winding $m{=}0$, the residual is purely symmetric: a spin-0,
colour-neutral scalar with three-fold internal structure.
Whether this corresponds to the physical Higgs boson requires
derivation of the mass generation mechanism, which is open.

%------------------------------------------------------------
\section{Confinement from Mixed Symmetry}
\label{sec:confinement}

The $n{=}2$, $m{=}1$ fermion (charged lepton) has a purely
antisymmetric residual.
Antisymmetry is self-contained: projection onto the antisymmetric
subspace recovers the residual exactly, with zero error.
The fermion can be described as a free asymptotic state.

The $n{=}3$, $m{=}1$ fermion (quark) has a mixed-symmetry residual.
Projection onto either self-contained class (symmetric or
antisymmetric) leaves a nonzero residual of magnitude
$\approx 0.408$, verified numerically.
A mixed-symmetry fermion always requires reference to the
configuration from which it arose.

\begin{conjecture}[Confinement from mixed symmetry]
A fermion whose compression residual $B$ has mixed symmetry cannot
be described as a free asymptotic state under any
minimum-description-length codec.
Colour confinement is a structural consequence of the three-site
codec geometry, not a separately imposed force law.
\end{conjecture}

%------------------------------------------------------------
\section{The Fermion Classification}
\label{sec:dictionary}

Table~\ref{tab:dictionary} presents the fermion complexity classes.
The right-hand column gives the Standard Model label for the
fermion event correlations that an observer using that language
would identify --- not the name of a fundamental particle, but
the name given to a pattern of fermion behaviour.

\begin{table}[h]
\centering
\begin{tabular}{cccclll}
\toprule
$(n,m)$ & $\|F\|$ & $\|B\|$ & Symmetry & Fermion class
        & SM correlation label & Status \\
\midrule
$(1,0)$ & $1$ & $0$ & — &
  Localised, no residual & vacuum & exact \\
$(2,0)$ & $1/\sqrt{2}$ & $1/\sqrt{2}$ & symmetric &
  Neutral lepton & neutrino & exact \\
$(2,1)$ & $1/\sqrt{2}$ & $1/\sqrt{2}$ & antisymmetric &
  Charged lepton & electron/muon/tau & exact \\
$(3,0)$ & $1/\sqrt{3}$ & $\sqrt{2/3}$ & symmetric &
  Colour-neutral scalar & Higgs sector & conjecture \\
$(3,1)$ & $1/\sqrt{3}$ & $\sqrt{2/3}$ & mixed &
  Colour-charged, confined & quark/gluon & exact \\
$(4,2)$ & $1/2$ & $\sqrt{3}/2$ & symmetric &
  Spin-2 fermion class & open & open \\
\bottomrule
\end{tabular}
\caption{Fermion complexity classes.
The ``SM correlation label'' column gives the Standard Model name
for the correlation pattern an observer would attribute to this
fermion class.
These are derived labels, not fundamental particles.
``Status'' refers to the $(n,m)$ norm structure of
Section~\ref{sec:theorems}, which is exact in every row; it does
\emph{not} certify that all SM particles sharing a label are
distinguished by the framework.
In particular the three charged leptons share the single class
$(2,1)$ and are not yet distinguished by mass --- see
Section~\ref{sec:masses}.}
\label{tab:dictionary}
\end{table}

%------------------------------------------------------------
\section{Mass and the Missing Generation Index}
\label{sec:masses}

The spectral complexity $C_s$ assigns an information cost
proportional to the dominant wavenumber of the fermionic wavefunction.
Under Solomonoff induction, $P(\psi) \propto 2^{-C_s(\psi)}$, so a
genuine mass-from-complexity relation would order observable fermions
by ascending $C_s$.

The $(n,m)$ classification of Section~\ref{sec:classes} does not yet
support this.
All three charged leptons --- electron, muon, tau --- occupy the
\emph{same} class, $(n,m) = (2,1)$.
The classification distinguishes charge and confinement
($m$ as winding number, $n$ as site count) but contains no third
parameter that could distinguish particles of different mass within
a class.
A genuine derivation of the lepton mass hierarchy from $C_s$ therefore
requires extending the classification with an explicit generation
index, which does not yet exist in this framework.

\begin{table}[h]
\centering
\begin{tabular}{lrrl}
\toprule
Fermion & Mass (MeV) & $\log_2(m/m_e)$ & Notes \\
\midrule
Electron &       $0.511$ &  $0.00$ & \\
Muon     &     $105.66$  &  $7.69$ & \\
Tau      &    $1776.86$  & $11.76$ & \\
\bottomrule
\end{tabular}
\caption{Charged lepton masses in $\log_2(m/m_e)$ space.
The two gaps are $7.69$ and $4.07$ bits respectively --- not equal
to each other, and not both close to a common integer unit.
No formula in this paper computes these numbers from $C_s$.}
\label{tab:masses}
\end{table}

A combinatorial argument is sometimes proposed as a candidate origin
for a discrete mass unit: one fermion hop requires four binary
choices (which site, which frame, real or imaginary component, sign
of phase), giving $2^4 = 16$ configurations and suggesting a 4-bit
granularity.
This argument is independently motivated but is \emph{not} connected
by any derivation to the $C_s$ formula of Paper~VI or to the observed
mass gaps in Table~\ref{tab:masses}.
No construction in this framework takes $(n,m) = (2,1)$ together with
a generation label and outputs $C_s = 8$ or $C_s = 12$.
The near-integer appearance of $7.69 \to 8$ relies on rounding a
number that is not itself derived, and the second gap ($4.07$) is
not close to the same unit that the first gap is rounded to.

We therefore do not claim a mass ladder.
The honest status is: the $(n,m)$ classification correctly separates
fermions by charge and confinement class, but says nothing yet about
why three charged-lepton masses exist within the single class
$(2,1)$, or what values they should take.
This is stated as an open problem in Section~\ref{sec:open} rather
than as a partial result.

%------------------------------------------------------------
\section{Discussion}
\label{sec:discussion}

\subsection{Why we only observe fermions}

The foundational asymmetry of the codec framework is this: fermions
are pixels, bosons are compression residuals.
A pixel can be observed --- it is a localised excitation at a
definite site, sampled by the Born rule.
A compression residual cannot be observed directly, because it is
not a configuration of the system --- it is the codec's internal
record of how the configuration changed.

Consider what actually happens in a particle physics experiment.
An electron enters a detector.
The detector registers a hit at a definite location.
Another electron, some distance away, is deflected.
Standard Model language inserts a photon between these two events
and says the photon ``mediated the interaction.''
But no detector ever registers a photon as a localised pixel.
What detectors register are always fermion events --- photoelectric
hits, track ionisations, Cherenkov rings.
The photon is the name we give to the correlation between those
events.

In the codec framework, that correlation is already present in the
off-diagonal of the density matrix.
The fermion in $n{=}2$, $m{=}1$ superposition has an antisymmetric
residual $B$ that encodes, exactly and completely, the correlation
structure between its transition events.
No separate photon particle is needed.
The Standard Model photon is a derived label, not a fundamental
object.

This is not a new claim in physics.
It echoes the Wheeler--Feynman absorber theory, the S-matrix
programme of the 1960s, and the modern amplitudes programme --- all
of which question whether the virtual particle picture is physically
necessary.
The codec framework makes the claim precise: virtual particles are
codec artifacts, and real particles are fermions.

\subsection{What the open problems reduce to}

The previous version of this paper listed eight open problems,
most of which concerned boson physics: W/Z mass generation, the
Higgs mechanism, gauge symmetry unification, the graviton
identification.
Under the fermion-first reframing, these dissolve or transform.

The W and Z bosons are observed as fermion recoil events at specific
energy scales.
The question is not ``what is the W boson?'' but ``why do fermion
transitions at those energy scales produce correlation patterns that
are massive and short-ranged?''
That is a question about the lognormal probability distribution of
Paper~V: fermion transitions near or above the lognormal peak are
expected to be suppressed and short-range.
This connection is plausible but not yet quantitative, and does not
rely on any claimed mass ladder.

The Higgs mechanism in the Standard Model gives mass to the W and Z
by coupling them to a scalar field.
In the codec framework, the prior question is more basic: why does
the $n{=}2$, $m{=}1$ fermion class contain particles of more than
one mass at all, when the $(n,m)$ classification of
Section~\ref{sec:classes} has no parameter to distinguish them
(Section~\ref{sec:masses})?
Until that generation index is identified, no statement about a
Higgs-like mass generation mechanism can be derived from this
framework; it remains speculative.

The graviton question is genuinely open and worth stating carefully.
Gravity was derived in Papers~IV--V at the geometric layer, from
the aspect ratio and counting equations.
The $n{=}4$, $m{=}2$ fermion class has codec spin 2.
Whether these are the same physical object requires a cross-layer
analysis.
We note the coincidence without claiming a derivation.

\subsection{The three genuine open problems}

Stripped of boson physics, the genuine open problems are:

\begin{enumerate}
  \item \textbf{What distinguishes the charged lepton generations.}
        The $(n,m)$ classification places electron, muon, and tau in
        the same class, $(2,1)$.
        No parameter in the present framework varies across them.
        Identifying such a parameter --- and only then asking whether
        it explains the observed masses, and whether it caps the
        count at three --- is the prerequisite open problem.

  \item \textbf{Why quarks have the masses they have.}
        The six quarks span several orders of magnitude and have not
        been connected to $C_s$ in any form, pending resolution of
        the generation-index problem above.

  \item \textbf{Why the fermion force laws have the form they have.}
        The $1/r^2$ force law was derived in Paper~VII from graviton
        flux conservation.
        The electromagnetic $1/r^2$ law should follow from the same
        argument applied to the $n{=}2$, $m{=}1$ fermion's
        antisymmetric residual.
        This derivation is not yet complete.
\end{enumerate}

Everything else --- boson masses, gauge symmetry groups, coupling
constants --- is a derived description of fermion correlation
patterns, not a fundamental open problem.

\subsection{The relationship between geometric and codec layers}

Papers~IV--V derived the geometric structure of spacetime from the
aspect ratio and counting equations, without reference to the
wavefunction codec.
Papers~VI--VIII derive quantum mechanics and particle structure from
the wavefunction codec, without reference to the geometric layer.

Both derivations are self-consistent.
The question is how they connect.
The natural bridge is the worldline: a fermion's trajectory through
spacetime is a complex worldline $\psi_{\mathrm{traj}}(t) = x(t) + iy(t)$,
and Paper~VII showed that the minimum spectral complexity closed
worldline is the ellipse --- Kepler's orbit.
The fermion codec and the geometric gravity are therefore not
independent: the minimum-$C_s$ worldline in the geometric spacetime
is exactly the trajectory selected by the fermion's Solomonoff
induction.
A full derivation of this bridge is the primary target of the
next paper.

%------------------------------------------------------------
\section{Open Problems}
\label{sec:open}

\begin{enumerate}

  \item \textbf{A generation index for charged leptons.}
        Identify a parameter, absent from the present $(n,m)$
        classification, that distinguishes electron, muon, and tau.
        Only once such a parameter exists can its connection (if any)
        to $C_s$, to the observed mass gaps, and to the number of
        generations be assessed.

  \item \textbf{Quark mass hierarchy.}
        Connect the six quark masses to $C_s$, pending resolution of
        the generation-index problem above.

  \item \textbf{Electromagnetic force law from fermion residual.}
        Apply the graviton flux argument of Paper~VII to the
        $n{=}2$, $m{=}1$ antisymmetric residual and derive the
        Coulomb $1/r^2$ law.

  \item \textbf{Confinement as theorem.}
        Prove that mixed-symmetry fermion residuals cannot be free
        asymptotic states under any minimum-description-length codec.

  \item \textbf{Cross-layer synthesis.}
        Connect the geometric gravity of Papers~IV--V to the
        fermion worldline codec.
        Determine whether the $n{=}4$, $m{=}2$ fermion class and
        the geometric graviton are the same physical entity.

  \item \textbf{Many-fermion extension.}
        Extend the single-fermion codec to tensor products,
        giving composite particles and the full Fock space.

\end{enumerate}

%------------------------------------------------------------
\section{Summary of Results}
\label{sec:results}

\textbf{Exact results:}
\begin{enumerate}
\item $\|F\| = 1/\sqrt{n}$ for any $n$-site equal superposition.
\item $\|B\| = \sqrt{(n-1)/n}$ for any $n$-site equal superposition.
\item $\|F\|^2 + \|B\|^2 = 1$ for any pure state.
\item Fermion spin-$1/2$ is encoded in the universal $-\pi/2$ phase
      of Paper~VII.
\item Pauli exclusion: $\|B\| = 0$ for any localised fermion.
\item Colour symmetry: $|B_{ij}| = 1/3$ for all colour pairs in the
      $n{=}3$, $m{=}1$ class.
\item Charge as winding number: $m = 0$ gives neutral fermion,
      $m = 1$ gives charged fermion, within the $n{=}2$ sector.
\end{enumerate}

\textbf{Conjectures (physically motivated, not yet derived):}
\begin{enumerate}
\item[8.] Confinement from mixed symmetry of the $n{=}3$, $m{=}1$
      fermion residual.
\item[9.] $n{=}3$, $m{=}0$ as the scalar sector underlying the
      Higgs mass generation pattern.
\end{enumerate}

\textbf{Open:}
\begin{enumerate}
\item[10.] A generation index distinguishing electron, muon, and tau
      within the single class $(2,1)$.
      No such parameter exists in the present framework; the
      previously claimed 4-bit mass ladder is not a derived
      consequence of $C_s$ and is withdrawn (see
      Section~\ref{sec:masses}).
\item[11.] Physical identification of the $n{=}4$, $m{=}2$ spin-2
      fermion class and its relation to geometric gravity.
\end{enumerate}

%------------------------------------------------------------
\section{Conclusion}
\label{sec:conclusion}

We have proved that the density matrix of a fermion in equal
superposition across $n$ sites splits orthogonally into a fermionic
diagonal $F$ and an off-diagonal residual $B$ satisfying
\[
  \|F\| = \frac{1}{\sqrt{n}}, \quad
  \|B\| = \sqrt{\frac{n-1}{n}}, \quad
  \|F\|^2 + \|B\|^2 = 1.
\]
The conservation law is exact for any pure state and subsumes the
Paper~VII result as a special case.

The central reframing of this paper is that bosons are not
fundamental particles.
They are the codec's internal record of fermion transitions.
Every physical observation is a fermion observation.
The Standard Model boson table is a derived description of fermion
correlation patterns, not an independent particle ontology.

Under this reframing, the particle classification is a
\emph{fermion} classification.
The pair $(n, m)$ describes the fermion's complexity class.
The $n{=}2$ sector gives charged leptons and neutrinos, distinguished
by winding number.
The $n{=}3$ sector gives confined colour-charged fermions (quarks)
and a scalar sector.
The classification does not yet distinguish the three charged lepton
generations, which all occupy the single class $(2,1)$; identifying
the missing generation index is left as an open problem rather than
asserted as a derived mass ladder.

The open problems reduce to three purely fermionic questions: what
distinguishes the lepton generations, why the quark masses have the
values they have, and whether the $1/r^2$ electromagnetic force law
follows from the fermion residual structure by the same argument
that gives gravity in Paper~VII.

%------------------------------------------------------------
\section*{Supplementary Material}

\begin{itemize}
\item \texttt{s1\_universal\_boson\_theorem.py} ---
        Theorem~\ref{thm:boson}: 1000 random phase trials per $n$.
\item \texttt{s2\_spin\_classification.py} ---
        Full $(n,m)$ classification; gluon colour structure.
\item \texttt{s3\_charge\_and\_confinement.py} ---
        Fractional occupancy; symmetry errors; confinement
        projection residuals.
\item \texttt{s4\_mass\_spectrum.py} ---
        Reports observed lepton $\log_2(m/m_e)$ values and the
        combinatorial 4-bit argument side by side, without asserting
        a connection between them; see Section~\ref{sec:masses}.
\item \texttt{s5\_fermion\_boson\_conservation.py} ---
        Theorems~\ref{thm:fermion} and \ref{thm:conservation};
        theta-dependent conservation; full spectrum summary.
\item \texttt{mass\_ladder\_audit.py} --- Audit the mass ladder conclusions.
\end{itemize}

\printbibliography

\end{document}
